
% 来源：EPUB part191.xhtml
\section{做胸无“大”志的人}

\BookQuestion{网友：为什么很多人起点和您一样，跑着跑着，差别越来越大呢？即使是同一方向也是望尘莫及呀。}

呵呵，可能是我比较笨，找不到捷径。（2010-03-27）

% 来源：EPUB part192.xhtml
\subsection{要脚踏实地}

我原话的意思是胸无“大”字（有时也戏称胸无大志）。这里的“大”是好大喜功的大。我说这话时流行“要做强先做大”，很多企业奔着世界500“大”去，所以那时候我就开始特别强调企业健康最重要，不要在乎“大”小，要着眼于脚踏实地。胸无“大”字和胸有大志不是矛盾的，实际上可以说它们是一回事。（2019-08-25）

\BookQuestion{网友：成功并能保持的人，都有他厉害的地方。}

很多所谓厉害的人其实仅仅是因为他们一直在老老实实做他们该做的事情而已。单看每件事情是很难看出来他们厉害在哪里的。（2018-11-05）

\BookQuestion{网友：什么算成功？}

成功的定义是什么很重要。我觉得能做自己喜欢做的事情就应该是一种成功，至于别人怎么看其实没那么重要。（2017-11-07）

\BookQuestion{网友：创业者总体是有渴望赢的野心的一群人。我直觉上认为这与胸无大志的平常心看似矛盾。}

有平常心的人们才更容易赢，才更有机会成功。胸无大志是相对于好大喜功说的。（2020-10-11）

\BookQuestion{网友：像我们智商情商普通、家庭资源关系普通、学历普通的青年人，如何突破自我和突破社会阶层，尽可能地快速实现自己的梦想和追求？}

我年轻时情况非常符合你说的那种青年人。不过我不太有“尽可能快速”的想法。（2020-12-03）

% 来源：EPUB part193.xhtml
\subsection{每天进步一点}

其实我自己对自己的定义就是普通人，别人怎么看我不太在乎。想起两句话：平常人很难有平常心，Common sense is not common（常识不常见）。（2011-03-29）

\BookQuestion{网友：您怎样理解不普通的人？}

不那么理解“不普通的人”到底是什么意思。我觉得绝大多数人都是普通人，但普通人通过积累是可以做出不普通的成就的，比如梅西。当然有些天赋会让人不那么普通，比如姚明的高度（姚明的成就实际上也是有巨大的努力在里面的），但那毕竟是极小比例的，也不是普通人能通过努力学到的东西。（2023-02-11）

大道确实是个普通人，高考虽然算是小考区状元，但进浙大时在班里35个人中考分排第18。考人大研究生时全班25个人考分排名第25。所以大道能有今天绝对不是因为绝顶聪明，而是因为一些其他的因素。比如，理性、想长远、不为清单……（2021-06-28）

\BookQuestion{网友：大道说自己是个普通人，但我一直不太理解。大道也许是个能力、智商比较普通的人，但不普通在于大道很理性，知道自己能做什么，然后只做自己能做的事，所以成功率要高很多。}

你的理解非常对！明白这点的人们20年后必有所获。（2021-06-29）

“我经常见到一些并不聪明的人成功，他们甚至也并不十分勤奋。但是他们都是一些热爱学习的‘学习机器’，他们每天晚上睡觉的时候都比早上起床时稍微多了那么一点点智慧。伙计，如果你前面有很长的路要走的话，这可是大有裨益的啊。”（芒格）

在坚持做对的事情的前提下，把事情做对的能力是可以学习的，只要每天进步一点就行。（2012-06-26）

哈，普通人一直尽量坚持做对的事情，几十年后就有机会封神了哈。（2024-10-30）
